\section{Worked Example 4: 安全评估和能力评估为什么会混在一起}

考虑一个能写 exploit 的模型。它有两种可能状态：
\begin{itemize}
    \item 能力很强，但默认拒答。
    \item 能力很强，而且愿意输出有害内容。
\end{itemize}

对 API 产品来说，前者和后者差异极大，因为用户实际上能看到的是 \textbf{propensity}。对开源模型来说，用户可以轻松改掉安全层，因此 \textbf{capability} 又变得更关键。

\begin{table}[H]
\centering
\begin{tabular}{p{3.0cm}p{4.2cm}p{5.8cm}}
\toprule
\textbf{视角} & \textbf{关心什么} & \textbf{为什么} \\
\midrule
API 产品 & propensity & 模型实际会不会输出有害内容 \\
开源模型 & capability & 安全策略能否被轻易移除 \\
监管机构 & 风险类别 & 是否符合政策、法律和部署要求 \\
\bottomrule
\end{tabular}
\caption{安全评估里 capability 和 propensity 是不同维度}
\end{table}

\begin{warningbox}{不要把“会做坏事”和“会不会说坏话”混为一谈}
有些评估在测模型的潜在能力，有些在测它在当前部署配置下的外显行为。两者都重要，但不能混淆。
\end{warningbox}

